% GENERATED FILE. Edit canonical lessons, then run npm run generate:books.
% canonical-source-hash: fnv1a64:09db6ea596f5b4bc
% canonical-lessons: BN-C03-ami, BN-C03-kemon, BN-C03-tumi-kemon-achho, BN-C03-bhalo, BN-C03-kono-bepar-na, BN-C03-practice

\chapter{How Are You?}
\label{ch:responding}
\hlchaptermodality{3}

\begin{chapteropening}
\textbf{By the end of this chapter:} \emph{I can ask how someone is in Bengali, answer naturally, and return the question.}

Complete BN-C03-practice, the canonical closing checkpoint, and demonstrate the chapter promise: ask how someone is in Bengali, answer naturally, and return the question.
\end{chapteropening}

\section[\bn{আমি}]{\bn{আমি} (āmi) — \textquotedblleft{}I,\textquotedblright{} cousin of English \emph{am}}

\label{lesson:BN-C03-ami}



\begin{quote}
To answer \textquotedblleft{}how are you?\textquotedblright{} you first need \textquotedblleft{}I\textquotedblright{} — and it hides the English word \emph{am}.
\end{quote}

\subsection*{The letters in this word}
\textbf{\bn{আ}} (independent long \emph{ā}) + \textbf{\bn{মি}} (\emph{mi}) $\to$ \textbf{\bn{আমি}} (\emph{āmi}).

\begin{cousinweb}
\textbf{\bn{আমি}} (\emph{āmi}, \textquotedblleft{}I\textquotedblright{}) descends from Sanskrit \textbf{\bn{অস্মি}} (\emph{asmi}, \textquotedblleft{}I am\textquotedblright{}), on the first-person stem \emph{asmad}. That Sanskrit \emph{asmi} is from Proto-Indo-European \textbf{*h\textsubscript{1}es-}, the \textquotedblleft{}to be\textquotedblright{} root — the very same word that gives English \textbf{am} (and Latin \emph{sum}). So Bengali's word for \textquotedblleft{}\textbf{I}\textquotedblright{} is, at the root, a form of \textquotedblleft{}\textbf{am}\textquotedblright{}: \textquotedblleft{}I\textquotedblright{} and \textquotedblleft{}am\textquotedblright{} were once one word. Its possessive is the \textbf{\bn{আমার}} (\emph{āmār}, \textquotedblleft{}my\textquotedblright{}) from Chapter 2.
\end{cousinweb}

\begin{grammarlens}[title={Grammar Lens: \textquotedblleft{}I\textquotedblright{} can often be dropped}]
Bengali verbs carry the person in their ending — \emph{āchhi} already means \textquotedblleft{}\textbf{I} am\textquotedblright{} — so \emph{āmi} is often left out once it's clear. You still learn it first, because the answer to \emph{kēmon āchho?} begins with it.
\end{grammarlens}

\subsection*{Your turn}
Say these aloud:

\begin{itemize}
\raggedright
  \item \textquotedblleft{}āmi\textquotedblright{}
  \item the pair — \emph{āmi} (I), \emph{āmār} (my)
  \item the English cousin hidden in it (\emph{am} — from the same \emph{asmi})
\end{itemize}

\begin{tcolorbox}[breakable,colback=teal!4,colframe=teal!35!black,title={Before you move on}]
What Sanskrit word is \emph{āmi} from, and its English cousin? (\emph{Asmi}, \textquotedblleft{}I am\textquotedblright{}; English \textbf{am} — same PIE root \emph{*h\textsubscript{1}es-}.)
\end{tcolorbox}
\section[\bn{কেমন}]{\bn{কেমন} (kēmon) — \textquotedblleft{}how,\textquotedblright{} another of the k- questions}

\label{lesson:BN-C03-kemon}



\begin{quote}
You met \textbf{\bn{কি}} (\emph{ki}, \textquotedblleft{}what\textquotedblright{}) in Chapter 2. Here is its sibling.
\end{quote}

\subsection*{The letters in this word}
\textbf{\bn{কে}} (\emph{kē}, \emph{kô} + the \emph{e}-sign) + \textbf{\bn{মন}} (\emph{mon}) $\to$ \textbf{\bn{কেমন}} (\emph{kēmon}).

\begin{cousinweb}
\textbf{\bn{কেমন}} (\emph{kēmon}, \textquotedblleft{}how, of what kind\textquotedblright{}) is from the interrogative stem \textbf{ka-}, from PIE \textbf{*k\textsuperscript{w}o-} — the same question-root behind English \textbf{wh-} and Latin \emph{qu-}. Bengali's questions nearly all begin with this \textbf{k-}: \emph{ki} (what), \emph{ke} (who), \emph{kôbe} (when), \emph{kōthāy} (where), \emph{kēmon} (how). It is the Indian branch of the same family of question-words that gave English \emph{what/who/how}.
\end{cousinweb}

\subsection*{Your turn}
Say these aloud:

\begin{itemize}
\raggedright
  \item \textquotedblleft{}kēmon\textquotedblright{}
  \item the k- question family — \emph{ki, ke, kôbe, kōthāy, kēmon}
\end{itemize}

\begin{tcolorbox}[breakable,colback=teal!4,colframe=teal!35!black,title={Before you move on}]
What root do Bengali's question-words share, and its English family? (The interrogative \emph{k-}, PIE \emph{*k\textsuperscript{w}o-}; the \emph{wh-} words.)
\end{tcolorbox}
\section[\bn{তুমি} \bn{কেমন} \bn{আছো}?]{\bn{তুমি} \bn{কেমন} \bn{আছো}? (tumi kēmon āchho?) — \textquotedblleft{}how are you?\textquotedblright{}}

\label{lesson:BN-C03-tumi-kemon-achho}



\begin{quote}
The everyday \textquotedblleft{}how are you?\textquotedblright{} — and the first time Bengali brings back a verb \textquotedblleft{}to be.\textquotedblright{}
\end{quote}

\subsection*{The letters in this word}
The new word is \textbf{\bn{আছো}} (\emph{āchho}, \textquotedblleft{}are you\textquotedblright{}), from the verb \textbf{\bn{আছা}} (\emph{āchhā}, \textquotedblleft{}to be, to exist, to have\textquotedblright{}) + the familiar-you ending \textbf{-\bn{ও}} (\emph{-o}). Note the \textbf{\bn{ছ}} (\emph{chh}) — a breathy \emph{ch}.

\begin{cousinweb}
\textbf{\bn{তুমি} \bn{কেমন} \bn{আছো}?} = \textbf{\bn{তুমি}} (\emph{tumi}, familiar \textquotedblleft{}you\textquotedblright{}) + \textbf{\bn{কেমন}} (\emph{kēmon}, \textquotedblleft{}how\textquotedblright{}) + \textbf{\bn{আছো}} (\emph{āchho}, \textquotedblleft{}are\textquotedblright{}) — \textquotedblleft{}you how are?\textquotedblright{}, the verb \textbf{last}.
\end{cousinweb}

\begin{grammarlens}[title={Grammar Lens: now the copula returns}]
In Chapter 2, Bengali dropped \textquotedblleft{}is\textquotedblright{}: \emph{āmār nām Arun} (\textquotedblleft{}my name Arun\textquotedblright{}). That \textbf{zero copula} works for \emph{equations} (X is Y). But \textquotedblleft{}how are you?\textquotedblright{} asks about a \textbf{state} — how you \emph{exist} — and for that Bengali uses a real verb, \textbf{\bn{আছা}} (\emph{āchhā}). So: no verb for \textquotedblleft{}my name is Arun,\textquotedblright{} a real verb for \textquotedblleft{}how are you.\textquotedblright{} Respectful form: \textbf{\bn{আপনি} \bn{কেমন} \bn{আছেন}?} (\emph{āpni kēmon āchhen?} — \emph{āchhen} for \emph{āpni}). Match the \textquotedblleft{}you\textquotedblright{} to the verb-ending.
\end{grammarlens}

\subsection*{Your turn}
Say these aloud:

\begin{itemize}
\raggedright
  \item the familiar question — \emph{tumi kēmon āchho?}
  \item the respectful version — \emph{āpni kēmon āchhen?}
  \item why here Bengali needs a verb but \textquotedblleft{}my name is Arun\textquotedblright{} didn't (state vs. equation)
\end{itemize}

\begin{tcolorbox}[breakable,colback=teal!4,colframe=teal!35!black,title={Before you move on}]
What verb does \textquotedblleft{}how are you?\textquotedblright{} use, and why does this need one when \textquotedblleft{}my name is Arun\textquotedblright{} didn't? (\emph{Āchhā}, \textquotedblleft{}to be/exist\textquotedblright{}; it describes a state, not an equation — the zero copula only covers X is Y.)
\end{tcolorbox}
\section[\bn{ভালো}]{\bn{ভালো} (bhālo) — \textquotedblleft{}well,\textquotedblright{} and how to answer}

\label{lesson:BN-C03-bhalo}



\begin{quote}
The word that answers \textquotedblleft{}how are you?\textquotedblright{} — and one Bengalis say all the time.
\end{quote}

\subsection*{The letters in this word}
\textbf{\bn{ভা}} (\emph{bhā}, an aspirated \emph{b} + long \emph{ā}) + \textbf{\bn{লো}} (\emph{lo}) $\to$ \textbf{\bn{ভালো}} (\emph{bhālo}).

\begin{cousinweb}
\textbf{\bn{ভালো}} (\emph{bhālo}, \textquotedblleft{}good, well, fine\textquotedblright{}) is the everyday Bengali \textquotedblleft{}good\textquotedblright{} — likely from Sanskrit \textbf{\bn{ভদ্র}} (\emph{bhadra}, \textquotedblleft{}good, auspicious, gentle\textquotedblright{}), worn down over centuries into \emph{bhālo}. It is one of the most-used words in the language: \emph{bhālo lāge} (\textquotedblleft{}{[}I{]} like it\textquotedblright{}), \emph{bhālobāsā} (\textquotedblleft{}love,\textquotedblright{} literally \textquotedblleft{}good-dwelling\textquotedblright{}).
\end{cousinweb}

\subsection*{The exchange — the reply}
With the \emph{āchhā} verb: \textbf{\bn{আমি} \bn{ভালো} \bn{আছি}} (\emph{āmi bhālo āchhi}) — \textquotedblleft{}\textbf{I am well}\textquotedblright{} (\emph{āchhi} = \textquotedblleft{}I am/exist\textquotedblright{}). The whole exchange:

\begin{quote}
— \emph{tumi kēmon āchho?} (\textquotedblleft{}How are you?\textquotedblright{}) — \emph{āmi bhālo āchhi, dhônnobad.} (\textquotedblleft{}I'm well, thank you.\textquotedblright{})
\end{quote}

\subsection*{Your turn}
Say these aloud:

\begin{itemize}
\raggedright
  \item \textquotedblleft{}bhālo\textquotedblright{}
  \item the full reply — \emph{āmi bhālo āchhi}
  \item the tender word built on it (\emph{bhālobāsā}, \textquotedblleft{}love\textquotedblright{})
\end{itemize}

\begin{tcolorbox}[breakable,colback=teal!4,colframe=teal!35!black,title={Before you move on}]
Give the full reply to \emph{tumi kēmon āchho?}, and name \emph{bhālo}'s likely Sanskrit source. (\emph{Āmi bhālo āchhi}; Sanskrit \emph{bhadra}, \textquotedblleft{}good.\textquotedblright{})
\end{tcolorbox}
\section[\bn{কোনো} \bn{ব্যাপার} \bn{না}]{\bn{কোনো} \bn{ব্যাপার} \bn{না} (kono bæpār nā) — \textquotedblleft{}no matter\textquotedblright{}}

\label{lesson:BN-C03-kono-bepar-na}



\begin{quote}
When someone thanks you — or apologises — this is the easy, gracious reply, and it reuses a word from Chapter 1.
\end{quote}

\subsection*{The letters in this word}
\textbf{\bn{কোনো}} (\emph{kono}, \textquotedblleft{}any\textquotedblright{}) + \textbf{\bn{ব্যাপার}} (\emph{bæpār}, \textquotedblleft{}matter, affair\textquotedblright{} — with the \emph{ya-phola} on \textbf{\bn{ব্য}}) + \textbf{\bn{না}} (\emph{nā}, \textquotedblleft{}no/not,\textquotedblright{} from Chapter 1) $\to$ \emph{kono bæpār nā}.

\begin{cousinweb}
\textbf{\bn{কোনো} \bn{ব্যাপার} \bn{না}} literally means \textquotedblleft{}\textbf{{[}it's{]} no matter}\textquotedblright{} — \emph{kono} (\textquotedblleft{}any\textquotedblright{}) + \emph{bæpār} (\textquotedblleft{}matter, affair\textquotedblright{}) + \emph{nā} (\textquotedblleft{}not\textquotedblright{}). It serves as \textquotedblleft{}it's nothing / no worries / never mind / you're welcome.\textquotedblright{} The word \textbf{\bn{ব্যাপার}} (\emph{bæpār}, \textquotedblleft{}matter, affair, business\textquotedblright{}) is \textbf{Sanskrit} — \emph{vyāpāra} (\textquotedblleft{}occupation, dealing\textquotedblright{}). And \textbf{\bn{না}} (\emph{nā}) you met in Chapter 1: the ancient negative, PIE \textbf{*ne}, behind English \textbf{no}, \textbf{not}, \textbf{none}. So the reply pairs a Sanskrit noun with the oldest negative in the family.
\end{cousinweb}

\subsection*{Your turn}
Say these aloud:

\begin{itemize}
\raggedright
  \item \textquotedblleft{}kono bæpār nā\textquotedblright{}
  \item what it literally means (\textquotedblleft{}{[}it's{]} no matter\textquotedblright{})
  \item the exchange — \emph{dhônnobad} / \emph{kono bæpār nā}
\end{itemize}

\begin{tcolorbox}[breakable,colback=teal!4,colframe=teal!35!black,title={Before you move on}]
What does \emph{kono bæpār nā} literally say, and which Chapter-1 word does it reuse? (\textquotedblleft{}It's no matter\textquotedblright{}; \emph{nā}, \textquotedblleft{}not\textquotedblright{} — PIE \emph{*ne}, English \emph{no}.)
\end{tcolorbox}
\section[Practice]{Chapter 3 — The how-are-you exchange}

\label{lesson:BN-C03-practice}



\begin{quote}
No new words. The whole responding cycle, from atoms you now own.
\end{quote}

\subsection*{The exchange}
Say these aloud:

\begin{itemize}
\raggedright
  \item \textquotedblleft{}How are you?\textquotedblright{} (familiar) — \emph{tumi kēmon āchho?}
  \item \textquotedblleft{}I'm well, thank you.\textquotedblright{} — \emph{āmi bhālo āchhi, dhônnobad.}
  \item \textquotedblleft{}No matter / you're welcome.\textquotedblright{} — \emph{kono bæpār nā.}
  \item the respectful question — \emph{āpni kēmon āchhen?}
\end{itemize}

\begin{cousinweb}
Say these aloud:

\begin{itemize}
\raggedright
  \item the k- question-word for \textquotedblleft{}how\textquotedblright{} (\emph{kēmon})
  \item \textquotedblleft{}I\textquotedblright{} and the verb \textquotedblleft{}to be/exist\textquotedblright{} (\emph{āmi}, \emph{āchhā})
  \item when Bengali needs the verb and when it drops it (state vs. equation)
\end{itemize}
\end{cousinweb}

\begin{grammarlens}[title={What you've built — roots you now carry}]
\begin{itemize}
\raggedright
  \item \emph{kēmon} $\leftarrow$ the interrogative \emph{k-} (PIE \emph{*k\textsuperscript{w}o-}), cousin of English \emph{how/what}
  \item \emph{āmi} $\leftarrow$ Sanskrit \emph{asmi}, cousin of English \textbf{am}
  \item \emph{bhālo} $\leftarrow$ Sanskrit \emph{bhadra} \textquotedblleft{}good\textquotedblright{}; \emph{āmi bhālo āchhi} \textquotedblleft{}I'm well\textquotedblright{}
  \item \emph{nā} $\leftarrow$ PIE \emph{*ne}, English \emph{no} (in \emph{kono bæpār nā})
\end{itemize}
\end{grammarlens}

\begin{tcolorbox}[breakable,colback=teal!4,colframe=teal!35!black,title={Before you move on}]
A stranger asks \emph{āpni kēmon āchhen?} — give a full, polite reply, and answer their thanks. (\emph{Āmi bhālo āchhi, dhônnobad} — and to thanks, \emph{kono bæpār nā}.)
\end{tcolorbox}
